\documentclass[11pt]{article}

\usepackage[utf8]{inputenc}
\usepackage[T1]{fontenc}
\usepackage{lmodern}
\usepackage{geometry}
\usepackage{amsmath,amssymb,amsfonts,amsthm,mathtools}
\usepackage{bm}
\usepackage{enumitem}
\usepackage{microtype}
\usepackage{hyperref}
\usepackage{titlesec}
\usepackage{booktabs}
\usepackage{array}
\usepackage{setspace}
\usepackage{mathrsfs}
\usepackage{physics}

\geometry{margin=1in}
\hypersetup{colorlinks=true, linkcolor=black, urlcolor=blue, citecolor=black}
\setlist[enumerate]{topsep=0.4em,itemsep=0.35em}
\setlist[itemize]{topsep=0.3em,itemsep=0.25em}
\titleformat{\section}{\large\bfseries}{\thesection.}{0.5em}{}
\titleformat{\subsection}{\normalsize\bfseries}{\thesubsection.}{0.5em}{}
\setstretch{1.06}

\newcommand{\Aether}{\AE{}ther}
\newcommand{\Aflow}{\AE{}ther-flow}
\newcommand{\TheoryName}{\Aflow{} Interpretation of Relativity}
\newcommand{\TheoryProgram}{\Aether{} / \Aflow{} framework}

\newcommand{\CompletedMetric}{g^{\sharp}}
\newcommand{\MatterSpace}{\Psi}

\newcommand{\Car}{\mathrm{car}}
\newcommand{\Cur}{\mathrm{cur}}
\newcommand{\Hom}{\mathrm{hom}}

\newcommand{\ForSpace}[1]{\mathcal I_{#1}^{\mathrm{for}}}
\newcommand{\PrimShell}{\mathcal S_{\mathrm{prim}}}
\newcommand{\Reduction}[1]{\mathcal R_{#1}}
\newcommand{\CurrentCarrier}[1]{\mathfrak C_{#1}^{\mathrm{cur}}}
\newcommand{\EqCarrier}[1]{\mathfrak C_{#1}^{\mathrm{eq}}}
\newcommand{\MetricOut}{\mathscr G_{\mathrm{out}}}
\newcommand{\MatterOut}{\mathscr T_{\mathrm{out}}}

\newcommand{\SubTarget}[1]{E_{#1}^{\mathrm{sub}}}
\newcommand{\EqnTarget}[1]{Q_{#1}^{\mathrm{eqn}}}
\newcommand{\HiddenSpace}[1]{\mathcal H_{#1}^{\mathrm{hid}}}
\newcommand{\HiddenBody}[1]{D_{#1}^{\mathrm{hid}}}

\newcommand{\SeedState}[1]{\mathcal S_{#1}^{\mathrm{seed}}}
\newcommand{\SeedBody}[1]{\mathcal B_{#1}^{\mathrm{seed}}}
\newcommand{\SeedShell}[1]{\mathcal Z_{#1}^{\mathrm{seed}}}
\newcommand{\SeedSupportShell}[1]{\mathcal U_{#1}^{\mathrm{seed,sup}}}
\newcommand{\SeedZeroCore}[1]{\mathcal Z_{#1}^{0,\mathrm{seed}}}
\newcommand{\SeedSplit}[1]{\Phi_{#1}^{\mathrm{seed}}}
\newcommand{\SeedCharge}[1]{\mathcal Q_{#1}^{\mathrm{seed}}}
\newcommand{\SeedEmbed}[1]{\zeta_{#1}^{\mathrm{seed}}}
\newcommand{\SeedKernel}[1]{\widetilde K_{#1}^{0,\mathrm{seed}}}
\newcommand{\SeedCoercive}[1]{P_{#1}^{\mathrm{seed,co}}}

\newcommand{\GermNucleus}{\mathfrak G^{\mathrm{sub,germ,zstn}}}
\newcommand{\GermState}[1]{\mathcal G_{#1}^{\mathrm{germ}}}
\newcommand{\GermBody}[1]{\mathcal B_{#1}^{\mathrm{germ}}}
\newcommand{\GermProj}[1]{\Pi_{#1}^{\mathrm{germ}}}
\newcommand{\GermSeedSec}[1]{\sigma_{#1}^{\mathrm{germ\leftarrow seed}}}
\newcommand{\GermSection}[1]{\Gamma_{#1}^{\mathrm{germ}}}
\newcommand{\GermShell}[1]{\mathcal Z_{#1}^{\mathrm{germ}}}

\newcommand{\ProtoPackage}{\mathfrak P^{\mathrm{sub,proto,nuc}}}
\newcommand{\ProtoState}[1]{\mathcal P_{#1}^{\mathrm{proto}}}
\newcommand{\ProtoBody}[1]{\mathcal B_{#1}^{\mathrm{proto}}}
\newcommand{\ProtoProj}[1]{\Pi_{#1}^{\mathrm{proto}}}
\newcommand{\ProtoGermMin}[1]{\mathfrak f_{#1}^{\mathrm{proto,germ}}}
\newcommand{\ProtoSection}[1]{\Gamma_{#1}^{\mathrm{proto}}}
\newcommand{\ProtoShell}[1]{\mathcal Z_{#1}^{\mathrm{proto}}}

\newtheorem{definition}{Definition}
\newtheorem{proposition}{Proposition}
\newtheorem{remark}{Remark}
\newtheorem{corollary}{Corollary}
\newtheorem{theorem}{Theorem}

\title{\textbf{The \TheoryName{}}\\[0.4em]
\large Primitive-Reservoir Observer-Reduction\\
\large Minimal Germ Zero-Shell Transport Nucleus\\
\large from Nucleus-Generating Substrate Proto-Germ Dynamics}
\author{}
\date{}

\begin{document}

\maketitle

\begin{abstract}
The current primitive-reservoir observer-reduction frontier is the minimal germ zero-shell transport nucleus. The previous collapse theorem already spent the remaining internal compression move available inside the observer-relevant germ seed-transport core: once the fixed downstream seed package is held in hand, the support-shell, zero-core, hidden-sector, realization, and solvability data are no longer independent benchmark-facing freedom beyond that smaller nucleus. After that result, the next honest benchmark-facing task is no longer another same-output relay through the former larger core. It is either to derive or justify the minimal germ zero-shell transport nucleus from more primitive explicit substrate structure, or to prove a genuinely smaller theorem-level collapse strictly inside that nucleus.

The present note takes the first route. For each sector it introduces \emph{nucleus-generating substrate proto-germ dynamics}: a compact convex proto-germ body over the recorded germ state space, a smooth proto-germ germ minimizer section, a smooth proto-germ restriction section, and an exact proto-germ shell projecting diffeomorphically to the exact germ shell. The current minimal germ zero-shell transport nucleus is then recovered canonically as the projected exact-shell transport image of that deeper proto-germ package. Because the already-recorded current frontier theorem reconstructs the omitted support-shell / zero-core / hidden-sector package from the nucleus together with the fixed seed target, the full observer-relevant germ seed-transport core follows downstream without any new larger datum.

The benchmark boundary remains fixed throughout. The one operative metric is still exactly \(\CompletedMetric\), the ordinary matter route is unchanged, there is no second operative metric, and the low-energy matter law is not enlarged. The positive result is therefore narrow but benchmark facing: the minimal germ zero-shell transport nucleus is no longer left as an unexplained primitive stopping point on the active line. The remaining honest burden moves one layer deeper again, to the proto-germ dynamics themselves or to a sharper collapse inside their observer-relevant core.
\end{abstract}

\section{Introduction}

The active derivational program of \TheoryName{} has now reached a cleaner burden than another same-output germ relay. The current germ-nucleus collapse theorem already proved that the observer-relevant germ seed-transport core contains no additional benchmark-facing freedom beyond a smaller minimal germ zero-shell transport nucleus. The accompanying routing consequence is immediate: the next honest benchmark-facing manuscript must either derive or justify that nucleus from more primitive substrate structure, or prove a genuinely smaller theorem inside it. Reopening the former larger germ seed-transport core, replaying the earlier full germ-package relay, or returning to older screened layers would no longer address the live burden.

The present note takes the first route. Its positive claim is not that final substrate microphysics has been identified, and it is not that the benchmark exact-relativistic package has been changed. Its claim is narrower and benchmark facing. The current minimal germ zero-shell transport nucleus is shown to arise canonically from a more primitive proto-germ-level substrate dynamics whose projected exact-shell transport already contains the live germ-nucleus data. In that sense the current nucleus is no longer the primitive place where the active line has to stop.

Two features matter here. First, the theorem targets the current live nucleus itself rather than revisiting the already-spent larger germ seed-transport core, so it does not spend the next move on an older same-output presentation. Second, the deeper proto-germ package is not merely another renaming of the current nucleus. It adds an explicit proto-germ state space, proto-germ body, proto-germ-to-germ projection, and deeper exact-shell transport data below the current frontier. The current germ body, seed anchoring, restriction section, and exact shell are therefore projected exact-shell images of a deeper substrate package rather than unexplained primitive data.

\paragraph{Framework and claim boundary.}
\TheoryProgram{} denotes the broader \Aether{} / \Aflow{} framework built on the claims that the \Aether{} is the underlying four-dimensional substrate of reality, the \Aflow{} is its intrinsic ordered motion, observed three-dimensional space is the local experiential slice of that deeper substrate, S-time is the experienced order of change arising from matter, light, and the \Aflow{}, and observed expansion is the three-dimensional appearance of deeper four-dimensional ordered motion. Within that framework, \TheoryName{} denotes the exact-closure theory in which the effective gravitational dynamics are adopted to be exactly Einsteinian with universal matter coupling. In the active sequence, \emph{adoption} means use of that established relativistic dynamics without claiming substrate derivation, while \emph{derivation} is reserved for a first-principles recovery from explicit substrate variables.

The present note stays strictly inside that boundary. It does not alter the benchmark exact-closure package. It does not introduce a second operative metric or a new low-energy matter law. It only proves that the current minimal germ zero-shell transport nucleus can itself be generated from a more primitive proto-germ-level substrate dynamics.

\section{Fixed Inputs}

\begin{definition}[Recorded primitive continuation data]
\label{def:fixed}
Fix the following continuation data from the active primitive-reservoir observer-reduction line.
\begin{enumerate}
    \item For each sector label \(\sigma\in\{\Car,\Cur,\Hom\}\), a primitive forcing shell
    \[
    \ForSpace{\sigma},
    \]
    a visible reduction map
    \[
    \Reduction{\sigma}:\ForSpace{\sigma}\to X_{\sigma},
    \]
    and shell-level current and equilibrium carriers
    \[
    \CurrentCarrier{\sigma}:\ForSpace{\sigma}\to Z_{\sigma}^{\mathrm{cur}},
    \qquad
    \EqCarrier{\sigma}:\ForSpace{\sigma}\to Z_{\sigma}^{\mathrm{eq}}.
    \]
    \item The product shell
    \[
    \PrimShell
    \equiv
    \ForSpace{\Car}\times \ForSpace{\Cur}\times \ForSpace{\Hom}.
    \]
    \item The recorded metric and ordinary-matter outputs
    \[
    \MetricOut:\PrimShell\to\{\text{observer metrics}\},
    \qquad
    \MatterOut:\PrimShell\times\MatterSpace\to\{\text{observer stress tensors}\},
    \]
    satisfying
    \begin{equation}
    \MetricOut(\mathbf w)=\CompletedMetric,
    \qquad
    \MatterOut(\mathbf w,\psi)
    =
    T_{\mu\nu}^{\mathrm{matter}}[\CompletedMetric,\psi]
    \label{eq:fixedoutputs}
    \end{equation}
    for every \(\mathbf w\in\PrimShell\) and every matter configuration \(\psi\in\MatterSpace\).
\end{enumerate}
\end{definition}

\begin{definition}[Recorded downstream seed package]
\label{def:seed}
Fix \(\sigma\in\{\Car,\Cur,\Hom\}\). The active line already fixes the downstream seed package at current benchmark-facing scope:
\begin{enumerate}
    \item a finite-dimensional seed state space \(\SeedState{\sigma}\), a compact convex seed body \(\SeedBody{\sigma}\subset\SeedState{\sigma}\), the exact seed shell \(\SeedShell{\sigma}\subset\SeedBody{\sigma}\), and the exact seed support shell \(\SeedSupportShell{\sigma}\subset\SeedShell{\sigma}\);
    \item a seed zero core \(\SeedZeroCore{\sigma}\subset\SeedSupportShell{\sigma}\), the same hidden fiber space \(\HiddenSpace{\sigma}\) with compact hidden body \(\HiddenBody{\sigma}\subset\HiddenSpace{\sigma}\), and the fixed seed split
    \[
    \SeedSplit{\sigma}:
    \SeedZeroCore{\sigma}\times\HiddenBody{\sigma}
    \xrightarrow{\cong}
    \SeedSupportShell{\sigma};
    \]
    \item the seed hidden charge
    \[
    \SeedCharge{\sigma}:\HiddenSpace{\sigma}\to\EqnTarget{\sigma},
    \]
    a smooth observer-compatible exact seed zero-response realization
    \[
    \SeedEmbed{\sigma}:\ForSpace{\sigma}\hookrightarrow\SeedZeroCore{\sigma},
    \]
    and the fixed-data solvability / coercive package recorded through \(\SeedKernel{\sigma}\) and \(\SeedCoercive{\sigma}\);
    \item benchmark faithfulness, meaning preservation of the benchmark outputs \eqref{eq:fixedoutputs}, no second operative metric, no enlarged low-energy matter law, and no premature promotion from adoption to completed derivation.
\end{enumerate}
\end{definition}

\begin{definition}[Recorded current minimal germ zero-shell transport nucleus]
\label{def:nucleus}
Fix \(\sigma\in\{\Car,\Cur,\Hom\}\). The recorded current minimal germ zero-shell transport nucleus \(\GermNucleus\) for sector \(\sigma\) consists of:
\begin{enumerate}
    \item the germ state space \(\GermState{\sigma}\), compact convex germ body \(\GermBody{\sigma}\subset\GermState{\sigma}\), projection
    \[
    \GermProj{\sigma}:\GermState{\sigma}\to\SeedState{\sigma},
    \]
    and seed section
    \[
    \GermSeedSec{\sigma}:\SeedBody{\sigma}\to\GermBody{\sigma}
    \]
    satisfying
    \[
    \GermProj{\sigma}\bigl(\GermBody{\sigma}\bigr)=\SeedBody{\sigma},
    \qquad
    \GermProj{\sigma}\circ\GermSeedSec{\sigma}
    =
    \mathrm{id}_{\SeedBody{\sigma}};
    \]
    \item the restriction section
    \[
    \GermSection{\sigma}:\GermState{\sigma}\to\SubTarget{\sigma},
    \]
    and the exact germ shell
    \[
    \GermShell{\sigma}
    \equiv
    \bigl(\GermSection{\sigma}\bigr)^{-1}(0)\cap\GermBody{\sigma},
    \]
    with
    \[
    \GermProj{\sigma}\big|_{\GermShell{\sigma}}
    :
    \GermShell{\sigma}
    \xrightarrow{\cong}
    \SeedShell{\sigma};
    \]
    \item benchmark faithfulness, meaning preservation of the benchmark outputs \eqref{eq:fixedoutputs}, no second operative metric, no enlarged low-energy matter law, and no premature promotion from adoption to completed derivation.
\end{enumerate}
By the already-recorded current frontier theorem, the support shell, zero core, hidden split, hidden charge, exact zero-response realization, and fixed-data solvability / coercive package are then canonically induced downstream from this nucleus together with the fixed seed package of Definition \ref{def:seed}, rather than taken as additional independent benchmark-facing data.
\end{definition}

\begin{remark}
Definitions \ref{def:fixed}, \ref{def:seed}, and \ref{def:nucleus} are fixed inputs. The primitive continuation data, the one operative metric, the ordinary matter route, the recorded downstream seed package, and the recorded current germ nucleus are not reopened here. The present note asks only whether the current minimal germ zero-shell transport nucleus can itself be generated from more primitive substrate dynamics.
\end{remark}

\begin{remark}[Downstream fixed fact]
\label{rem:downstream}
The already-recorded current frontier theorem proves that once a minimal germ zero-shell transport nucleus is fixed, the larger observer-relevant germ seed-transport core is canonically recovered and carries no extra benchmark-relevant freedom. The present note uses that downstream fact only as a routing consequence: after the induced nucleus is obtained, no second larger current-core package needs to be postulated anew.
\end{remark}

\section{Nucleus-Generating Substrate Proto-Germ Dynamics}

\begin{definition}[Nucleus-generating substrate proto-germ dynamics]
\label{def:proto}
Fix \(\sigma\in\{\Car,\Cur,\Hom\}\). A nucleus-generating substrate proto-germ dynamics package \(\ProtoPackage\) for sector \(\sigma\) consists of the following data.
\begin{enumerate}
    \item A finite-dimensional Euclidean proto-germ state space \(\ProtoState{\sigma}\), a compact convex proto-germ body
    \[
    \ProtoBody{\sigma}\subset\ProtoState{\sigma}
    \]
    with nonempty interior, an affine surjection
    \[
    \ProtoProj{\sigma}:\ProtoState{\sigma}\to\GermState{\sigma},
    \]
    and a smooth proto-germ fiber minimizer
    \[
    \ProtoGermMin{\sigma}:\GermBody{\sigma}\to\ProtoBody{\sigma}
    \]
    satisfying
    \[
    \ProtoProj{\sigma}\bigl(\ProtoBody{\sigma}\bigr)=\GermBody{\sigma},
    \qquad
    \ProtoProj{\sigma}\circ\ProtoGermMin{\sigma}
    =
    \mathrm{id}_{\GermBody{\sigma}}.
    \]
    \item A smooth proto-germ restriction section
    \[
    \ProtoSection{\sigma}:\ProtoState{\sigma}\to\SubTarget{\sigma},
    \]
    and the exact proto-germ shell
    \[
    \ProtoShell{\sigma}
    \equiv
    \bigl(\ProtoSection{\sigma}\bigr)^{-1}(0)\cap\ProtoBody{\sigma},
    \]
    satisfying
    \[
    \ProtoSection{\sigma}\circ\ProtoGermMin{\sigma}
    =
    \GermSection{\sigma},
    \qquad
    \ProtoProj{\sigma}\big|_{\ProtoShell{\sigma}}
    :
    \ProtoShell{\sigma}
    \xrightarrow{\cong}
    \GermShell{\sigma}.
    \]
\end{enumerate}
The package is \emph{benchmark faithful} if it preserves the benchmark outputs \eqref{eq:fixedoutputs}, introduces no second operative metric, introduces no enlarged low-energy matter law, and is presented as a deeper substrate-side source for the current minimal germ zero-shell transport nucleus rather than as a replacement benchmark theory.
\end{definition}

\begin{remark}
Definition \ref{def:proto} is intentionally narrower than another full germ relay. It does not rebuild the former observer-relevant germ seed-transport core as an independent datum, and it does not spend the next move on another restatement of the earlier full germ-generating substrate proto-germ package. It only introduces the deeper state space, body, and exact-shell transport data needed to generate the current nucleus.
\end{remark}

\begin{remark}
The label \emph{proto-germ} is used here only as a local marker for a deeper layer below the current germ-nucleus frontier. It does not by itself assert a completed microscopic ontology or a final primitive law. Its role is only to name the next sharper transport package below the current nucleus.
\end{remark}

\section{Projected Exact-Shell Transport to the Current Nucleus}

\begin{proposition}[The seed anchor persists one layer deeper]
\label{prop:anchor}
Fix a benchmark-faithful proto-germ package. Then
\[
(\GermProj{\sigma}\circ\ProtoProj{\sigma})\bigl(\ProtoBody{\sigma}\bigr)
=
\SeedBody{\sigma},
\]
\[
(\GermProj{\sigma}\circ\ProtoProj{\sigma})\circ
\ProtoGermMin{\sigma}\circ\GermSeedSec{\sigma}
=
\mathrm{id}_{\SeedBody{\sigma}},
\]
and the composite shell transport
\[
(\GermProj{\sigma}\circ\ProtoProj{\sigma})\big|_{\ProtoShell{\sigma}}
:
\ProtoShell{\sigma}
\xrightarrow{\cong}
\SeedShell{\sigma}
\]
is a diffeomorphism.
\end{proposition}

\begin{proof}
Definition \ref{def:proto}(1) gives
\[
\ProtoProj{\sigma}\bigl(\ProtoBody{\sigma}\bigr)=\GermBody{\sigma}
\]
and
\[
\ProtoProj{\sigma}\circ\ProtoGermMin{\sigma}
=
\mathrm{id}_{\GermBody{\sigma}}.
\]
Applying Definition \ref{def:nucleus}(1) to those identities yields the first two displayed formulas.

For the shell transport, Definition \ref{def:proto}(2) gives a diffeomorphism
\[
\ProtoProj{\sigma}\big|_{\ProtoShell{\sigma}}
:
\ProtoShell{\sigma}
\xrightarrow{\cong}
\GermShell{\sigma},
\]
while Definition \ref{def:nucleus}(2) gives a diffeomorphism
\[
\GermProj{\sigma}\big|_{\GermShell{\sigma}}
:
\GermShell{\sigma}
\xrightarrow{\cong}
\SeedShell{\sigma}.
\]
Their composition is therefore a diffeomorphism from \(\ProtoShell{\sigma}\) to \(\SeedShell{\sigma}\).
\end{proof}

\begin{proposition}[Proto-germ dynamics canonically reproduce the current germ nucleus]
\label{prop:nucleus}
Fix a benchmark-faithful proto-germ package. Then projected exact-shell transport along
\[
\ProtoProj{\sigma}:\ProtoState{\sigma}\to\GermState{\sigma}
\]
reproduces exactly the recorded current minimal germ zero-shell transport nucleus of Definition \ref{def:nucleus}.
\end{proposition}

\begin{proof}
The current nucleus body and its seed-anchored sectioning data are recovered by Definition \ref{def:proto}(1): the proto-germ body projects onto the recorded germ body, the proto-germ minimizer sections that projection on \(\GermBody{\sigma}\), and Proposition \ref{prop:anchor} shows that the fixed seed anchor \(\GermProj{\sigma},\GermSeedSec{\sigma}\) persists unchanged one layer deeper.

The current nucleus restriction data are recovered by Definition \ref{def:proto}(2): the recorded germ section is the minimizing transport of the proto-germ section, and the exact current shell is the projected image of the exact proto-germ shell under a diffeomorphism.

Those are precisely the constituents retained by the current minimal germ zero-shell transport nucleus frontier. Therefore the recorded current nucleus is not postulated independently inside this deeper class. It is the projected exact-shell transport image of the proto-germ package.
\end{proof}

\begin{proposition}[The benchmark package remains fixed]
\label{prop:boundary}
For every benchmark-faithful proto-germ package, the operative metric remains exactly \(\CompletedMetric\), the ordinary matter route remains unchanged, no second operative metric is introduced, and the low-energy matter law is not enlarged.
\end{proposition}

\begin{proof}
This is part of the benchmark-faithfulness requirement in Definition \ref{def:proto}. Equation \eqref{eq:fixedoutputs} remains unchanged on the full primitive forcing shell, so the induced germ nucleus does not alter the benchmark exact-closure package.
\end{proof}

\begin{proposition}[The larger current core is still recovered downstream]
\label{prop:downstream}
Fix a benchmark-faithful proto-germ package. Once projected exact-shell transport reproduces the current minimal germ zero-shell transport nucleus, the already-recorded current frontier theorem canonically recovers the full observer-relevant germ seed-transport core without any new larger datum.
\end{proposition}

\begin{proof}
Proposition \ref{prop:nucleus} recovers the recorded current nucleus. By Remark \ref{rem:downstream}, the already-recorded current frontier theorem then applies verbatim and canonically supplies the support shell, zero core, hidden split, hidden charge, exact zero-response realization, and fixed-data solvability / coercive package from that nucleus together with the fixed seed package. No second independently chosen larger current-core datum is added here.
\end{proof}

\section{Main Theorem}

\begin{theorem}[Minimal germ zero-shell transport nucleus from nucleus-generating substrate proto-germ dynamics]
\label{thm:main}
Fix the primitive continuation data of Definition \ref{def:fixed}, the recorded downstream seed package of Definition \ref{def:seed}, and the recorded current germ nucleus of Definition \ref{def:nucleus}. Then, for each sector \(\sigma\in\{\Car,\Cur,\Hom\}\), every benchmark-faithful nucleus-generating substrate proto-germ dynamics package canonically reproduces the current minimal germ zero-shell transport nucleus with the same benchmark one-metric outputs. More precisely:
\begin{enumerate}
    \item projected exact-shell transport along \(\ProtoProj{\sigma}\) reproduces exactly the recorded current nucleus by Proposition \ref{prop:nucleus};
    \item the current minimal germ zero-shell transport nucleus is therefore no longer an unexplained primitive stopping point on the active line, but the projected exact-shell transport image of a more primitive proto-germ dynamics;
    \item the benchmark boundary remains fixed by Proposition \ref{prop:boundary};
    \item by Proposition \ref{prop:downstream}, the full observer-relevant germ seed-transport core is then recovered downstream from that induced nucleus without any new larger datum;
    \item consequently the remaining honest burden moves one layer deeper again: derive or justify the nucleus-generating substrate proto-germ dynamics themselves from still more primitive structure, or else prove a genuinely smaller theorem-level collapse, sharper necessity result, or sharper uniqueness result strictly inside the proto-germ package.
\end{enumerate}
\end{theorem}

\begin{proof}
Item (1) is Proposition \ref{prop:nucleus}. Item (2) is the project-level consequence of item (1): once the proto-germ package is fixed, the current minimal germ zero-shell transport nucleus is no longer introduced as an independent primitive tuple. It is the projected exact-shell transport image of a deeper state space, deeper body, deeper section, and deeper exact shell.

Item (3) is Proposition \ref{prop:boundary}. Item (4) is Proposition \ref{prop:downstream}. Item (5) is the routing consequence of the previous items together with the active safeguard against counting same-output deeper relays as new front-line progress once the live observer-relevant datum is unchanged.
\end{proof}

\begin{corollary}[What must be done next]
\label{cor:next}
After Theorem \ref{thm:main}, the next honest benchmark-facing manuscript must either:
\begin{enumerate}
    \item derive or justify the nucleus-generating substrate proto-germ dynamics from still more primitive substrate structure in a way that changes benchmark-facing status;
    \item identify a genuinely smaller theorem-level observer-relevant core strictly inside the proto-germ package and prove a sharper collapse to it;
    \item do either of the previous two while preserving the same benchmark one-metric package and the same claim boundary.
\end{enumerate}
Another same-output restatement of the current minimal germ zero-shell transport nucleus, another same-output deeper-origin relay that leaves that nucleus unchanged, another reopening of the former observer-relevant germ seed-transport core as if it were still the live burden, another replay of the earlier full germ-package handoff, a reopening of older seed, root, foundation, ground, origin, precursor, lift, shell-restriction, selector, spectral, or stationary layers, or any move that introduces a second operative metric, enlarges the ordinary matter law, or uses stronger promotion language does not satisfy the present burden.
\end{corollary}

\begin{proof}
Theorem \ref{thm:main} removes the current germ nucleus as the unexplained stopping point on the active line. Therefore another manuscript that merely restates that same nucleus, or merely returns to an already-screened larger relay, does not advance the benchmark-facing status of the project. What remains open is either to go one honest layer deeper again, to the proto-germ dynamics themselves, or to discover a smaller theorem-level observer-relevant core inside the proto-germ package.
\end{proof}

\section{Conclusion}

The positive result of this note is narrow but benchmark facing. The current germ-nucleus collapse theorem had already shown that the larger observer-relevant germ seed-transport core contains no extra benchmark-relevant freedom beyond a smaller minimal germ zero-shell transport nucleus. The present theorem then removes the next unexplained stopping point by showing that this nucleus itself is the projected exact-shell transport image of a more primitive nucleus-generating substrate proto-germ dynamics package.

The scope remains disciplined. The benchmark package is unchanged: the one operative metric is still exactly \(\CompletedMetric\), the ordinary matter route is unchanged, there is no second operative metric, and the low-energy matter law is not enlarged. Nothing here upgrades adoption to completed derivation or licenses stronger promotion language.

The open burden therefore moves one honest step further again. The next benchmark-facing manuscript should derive or justify the proto-germ dynamics from still more primitive substrate structure, unless the sharper honest gain available instead is a genuinely smaller theorem-level collapse, sharper necessity result, or sharper uniqueness result inside the proto-germ package. Until then, the current result should be read for what it is: a deeper germ-nucleus derivation on the active line of \TheoryName{}, not a claim that the full first-principles program is finished.

\end{document}
